

In \cref{constr_sync:sec:pair_aut},
we generalize the pair criterion (\cref{prelim:thm:pair_sync_condition})
to constraint languages that are monoids.
Then, in \cref{constr_sync:monoid-lang:shortest-sync-word}, we derive algorithmic consequences for testing constrained synchronizability and for the computation of synchronizing words of at most cubic length, as well as sufficient conditions for containment in \NL.

\subsection{A Pair Criterion}
\label{constr_sync:sec:pair_aut}

\Cref{constr_sync:thm:pair_condition} generalizes the pair criterion (\cref{prelim:thm:pair_sync_condition})
to the case of monoid constraint languages.
Further, in \cref{constr_sync:prop:pair_bound}, utilizing the pair automaton,
we give a quadratic upper bound for shortest words in any constraint language, even those that are not monoids,
mapping two states to a common state (we also give time and space bounds for computing such a word).

\begin{theorem} 
	\label{constr_sync:thm:pair_condition}
	Suppose $L \subseteq \Sigma^*$ is a monoid language.
	An automaton $\semiautomat{A}{Q}{\Sigma}{\delta}$
	has a synchronizing word in $L$ if and only if
	for all $p, q \in Q$ there exists a word $u \in L$
	such that $\applytrans{\delta}{p}{u} = \applytrans{\delta}{q}{u}$.  
\end{theorem}
\begin{proof}
	Set $n = |Q|$. A synchronizing word in $L$
	maps all pairs of states to the same state.
    Conversely, assume all pairs of states can be mapped to a single state
	by a word in~$L$. 
	Then, set $Q_1 = Q$ and, for $i \in \{1,\ldots,n-1\}$,
	$Q_{i+1} = \applytrans{\delta}{Q_i}{u_i}$
	where, 	if $|Q_i| > 1$, let $u_i \in L$ be
	such that $\applytrans{\delta}{p}{u_i} = \applytrans{\delta}{q}{u_i}$
	for two distinct $p,q \in Q_i$, and $u_i = \varepsilon$ otherwise.
	Since $|Q_{i+1}| < |Q_i|$ for $|Q_i| > 1$ and $|Q_1| = n$, 
	there exists $j \in \{1,\ldots,n-1\}$ such that $|Q_j|=1$.
	Set $u = u_1 \cdots u_j$.
	We have $u \in \Syn(\mathcal A)$.
	Since $L$ is a monoid we get $u \in L$.
\end{proof}

\begin{remark}
	We opted to present a direct proof for \cref{constr_sync:thm:pair_condition} above.
	Another proof along the lines as in the unconstrained case presented in \cref{prelim:sync:pair}
	is possible
	by (1) first  showing that a trap-state automaton
	has a synchronizing word in a monoid language~$L$
	if and only if there exists for each state
	a word in~$L$ mapping this state to the trap state
	and then (2) using
	that an arbitrary automaton has the same
	set of synchronizing words as the pair automaton
	in which the singleton sets are merged into a trap state.	
\end{remark}

\begin{proposition}
	\label{constr_sync:prop:pair_bound}
	Suppose $\automat{B}{P}{\Sigma}{\mu}{p_0}{F}$ is a PDFA.
	Let $\semiautomat{A}{Q}{\Sigma}{\delta}$ be an $n$-state automaton
	and $p, q \in Q$. In time $O(|P|n^2 + |\Sigma||P|n^2)$
	and working space $O(|\Sigma||P|n^2)$
	we can test whether there exists $u \in L(\mathcal B)$ such that $\applytrans{\delta}{p}{u} = \applytrans{\delta}{q}{u}$
	and, if this is the case,
	compute a shortest $u$, where we have
	\[
	|u| < \binom{n}{2} \cdot |P| + |P|.
	\]
\end{proposition} 
\begin{proof}
	Consider the product automaton $( \mathcal A^{\le 2} / {\sim_S} ) \times \mathcal B$
	where $S = \{ \{q\} \mid q \in Q \}$.
	A shortest word in~$u \in L(\mathcal B)$ mapping two states $p,q \in Q$
	to a single state corresponds
	to a path from 
	$(\{p,q\}, p_0)$ to a state $(S, r)$ with $r \in F$.
	Since $( \mathcal A^{\le 2} / {\sim_S} ) \times \mathcal B$
	has  $(\binom{n}{2} + 1) \cdot |P|$ states,
	we obtain $|u| < \binom{n}{2} \cdot |P| + |P|$.
	Now, $( \mathcal A^{\le 2} / {\sim_S} ) \times \mathcal B$
    with $|P| \cdot (\binom{n}{2} + 1)$ states
	and at most $|\Sigma||P| (\binom{n}{2} + 1)$ transitions
	can be computed in 	time and space $O(|\Sigma||P|n^2)$.   
	Then, existence of $u$, and its computation in case it exists, can be done by 
	a breadth-first search~\cite{cormen01introduction}
	in  time $O(|P|n^2 + |\Sigma||P|n^2)$
	and space $O(|P|n^2)$.
\end{proof}

\subsection{Computing Synchronizing Words and Containment in~\texorpdfstring{$\NL$}{NL}}
\label{constr_sync:monoid-lang:shortest-sync-word}

\Cref{constr_sync:thm:monoid_compute_sync_word} 
gives a cubic bound on shortest synchronizing words
in constraint languages that are monoids
and time and space upper bound to compute such a word.
The algorithm is based on the method proofs used for
\cref{constr_sync:thm:pair_condition,constr_sync:prop:pair_bound}.
Then, as a consequence of \cref{constr_sync:thm:pair_condition,constr_sync:prop:pair_bound},
in \cref{thm:reg_monoid} we conclude that for monoid constraint languages, the constrained synchronization problem is in \NL.

\begin{theorem}\label{constr_sync:thm:monoid_compute_sync_word}
	Let $\automat{B}{P}{\Sigma}{\mu}{p_0}{F}$
	be a PDFA
	and $L = L(\mathcal B)$ be a monoid.
	Suppose $\semiautomat{A}{Q}{\Sigma}{\delta}$
	is an $n$-state automaton.	
	In time $O(|\Sigma||P|n^2)$
	we can test whether $\mathcal A$
	has a synchronizing word in~$L$.
	Testing whether~$\mathcal A$ has a synchronizing word
	and, if this is the case,
	computing one of length at most 
	$(n - 1) \cdot (  \binom{n}{2} \cdot |P| + |P| - 1 )$
	can be done in time $O(|\Sigma||P|n^3)$ and working space $O(|\Sigma||P|n^2)$ (not counting the space the output needs).
\end{theorem}
\begin{proof}
	To test whether the automaton $\mathcal A$ has a 
	synchronizing word in $L(\mathcal B)$, we construct 
	$( \mathcal A^{\le 2} / {\sim_S} ) \times \mathcal B$ and check whether, 
	from every state $({p,q}, p_0)$ with $p,q \in Q$, a state in ${S} \times F$ is reachable.
	This can be done using a breadth-first search that runs ``backward''  starting from the states in ${S} \times F$.
	The algorithm runs in time $O(|\Sigma||P|n^2)$ and is correct by \cref{constr_sync:thm:pair_condition}.
		
	Next, to give an algorithms that also computes a synchronizing word, 
	first consider the proof of \cref{constr_sync:thm:pair_condition}.
	We can compute the sets $Q_i$ using a loop that iterates at most $n - 1$ times.	
	By \cref{constr_sync:prop:pair_bound}, in each iteration, for $|Q_i| > 1$,
	we can test whether a word $u_i$
	with $|\delta(Q_i, u_i)| < |Q_i|$
	exists in time 
	$O(|P|n^2 + |\Sigma||P|n^2)$
	and working space $O(|\Sigma||P|n^2)$,
	and, moreover, compute a shortest $u_i$ where 
	\[
	|u_i| \le \binom{n}{2} \cdot |P| + |P| - 1.
	\]
	If there exists no such word in some iteration, then, by \cref{constr_sync:thm:pair_condition},
	there exists no synchronizing word in $L$.
	Otherwise, concatenating the $u_i$
	gives a synchronizing word in~$L$
	of length at most $(n - 1) \cdot (  \binom{n}{2} \cdot |P| + |P| - 1 )$.
\end{proof}

The algorithm given in the proof of \cref{constr_sync:thm:monoid_compute_sync_word} is an adaption 
of a similar algorithm for the unconstrained case~\cite{Eppstein90,Sandberg04,DBLP:conf/lata/Volkov08}.
 
\begin{theorem}
	\label{thm:reg_monoid}
	Let $L \subseteq \Sigma^*$ be a regular language that is a monoid.\footnote{We avoided to call 
		it a regular monoid, as the notion of regularity is also used in semigroup theory with a different meaning (see, for example,~\cite{howie1995fundamentals}).}
	Then $\constrsync{L}$
	is in~\NL. 
\end{theorem} 
\begin{proof} 
	Suppose $\automat{B}{P}{\Sigma}{\mu}{p_0}{F}$ is a PDFA with $L = L(\mathcal B)$.
	Let $\semiautomat{A}{Q}{\Sigma}{\delta}$
	and $n = |Q|$.
	By \cref{constr_sync:thm:pair_condition,constr_sync:prop:pair_bound},
	the problem can be decided in \NL\ 
	by cycling through all pairs of states $p,q \in Q$
	and nondeterministically guess letter-by-letter
	$u \in \Sigma^*$ with $|u| \le |P|\cdot ( \binom{n}{2} + 1 )$, $u \in L$,
	and $\delta(p, u) = \delta(q, u)$.
	If for all state pairs an appropriate word is found,
	the automaton admits a synchronizing word in $L$, if not, we reject.
\end{proof} 

Combining Theorem~\ref{thm:reg_monoid} and~\ref{constr_sync:thm:regular_prefix_codes_and_automata},
we get the next corollary.

\begin{corollary}
	\label{constr_syn:cor:start_equals_final_state}
	For $\automat{B}{P}{\Sigma}{\mu}{p_0}{\{p_0\}}$,
	$\constrsync{L(\mathcal B)}$ is in~\NL.
\end{corollary}

Next, we derive a sufficient condition on the concatenation with monoid
languages resulting in problems in \NL.

\begin{corollary}[UVW-Condition]
	\label{constr_sync:cor:UVW}
	Let $U,V, W \subseteq \Sigma^*$ be such that 
	$U \subseteq \suff{V}$, $W \subseteq \pref{V}$
	and $V$ is a regular language that is a monoid.
	Then 
	$
	\constrsync{UVW} \eqACZEROm \constrsync{V}
	$
	and $\constrsync{UVW}$ is in \NL.
\end{corollary}
\begin{proof}
	As $V = V^*$, the assumptions give $\factor{UVW} = \factor{V}$. Then 
	Corollary~\ref{gen_results:cor:factors_etc_complexity} and Theorem~\ref{thm:reg_monoid} yield the claim.
\end{proof}

The property that $V$ is a monoid is essential in \cref{constr_sync:cor:UVW}.
For example, as listed in Theorem~\ref{constr_sync:thm:introductory_examples},
$\constrsync{ab^*} \in \LOGSPACE$,
but $\constrsync{ab^*a}$ is \NP-complete.

Lastly, by \cref{constrsync:thm:monoid-sparseregular-or-exp}, a sparse language that is a monoid is either
an infinite subset of $u^*$ for some $u \in \Sigma^*$
or equals $\{\varepsilon\}$. The language $\{\varepsilon\}$ is the only finite monoid language (and so, for this language
an input semiautoamton has a synchronizing word if and only if it has a single state).
Thus, the constrained synchronization problem with such a constraint language is either solvable in constant
time or \LOGSPACE-complete by \cref{constr_sync:thm:infinite,constr_sync:cor:in_star_star}.

\subsection{Strongly Connected Constraint Automata}
\label{sec:constr_sync:SCC} 



\Cref{constrsync:thm:SCC-prefix-code} yields with \cref{constr_sync:thm:regular_prefix_codes_and_automata}
that the the computational complexities arising from 
strongly connected constraint PDFAs equal those arising
from constraint languages of the form $C^*$ (which are monoids)
for $C=\emptyset$ or a prefix code $C$ that is a regular language. (Recall that these
languages equal those of the form $L(\mathcal B_{p_0, \{p_0\}}$ by \cref{constr_sync:thm:regular_prefix_codes_and_automata}.)
\Cref{constrsync:cor:SCC} shows that the constrained synchronization problem is always in \NL\ 
for strongly connected constraint automata. As a special case, in
\cref{constr_sync:cor:one-state} we classify the complexity for one-state constraint automata.

Observe that for a strongly connected PDFA $\automat{B}{P}{\Sigma}{\mu}{p_0}{F}$, $L(\mathcal B)$ is finite if and only
if $L(\mathcal B) = \emptyset$ (equivalently $F = \emptyset$)
or $L(\mathcal B) = \{\varepsilon\}$ (equivalently $P = \{p_0\}$ and we have no 
outgoing transitions at $p_0$).

\begin{theorem} \label{constrsync:thm:SCC-prefix-code}
	Let $\automat{B}{P}{\Sigma}{\mu}{p_0}{F}$ be a strongly connected PDFA with $F \ne \emptyset$.
	Then $\constrsync{L(\mathcal B)} \eqACZEROm \constrsync{C^*}$
	where $C$ is some prefix code $C \subseteq \Sigma^+$
	that is regular or we have  $C = \emptyset$.
\end{theorem}
\begin{proof}
	A strongly connected PDFA is trim and every state is accessible
	and coaccessible. Thus, by \cref{gen_results:cor:factors_etc_complexity,constr_sync:prop:final_states}, 
	we find\footnote{For language of the form $L(\mathcal B_{p_0, \{p_0\}})$
		by discarding useless states, i.e., assming~$\mathcal B$ to be trim,
		we can assume that $\mathcal B$ is strongly connected.}
	\[
	\constrsync{L(\mathcal B)} \eqACZEROm \constrsync{L(\mathcal B_{p_0, \{p_0\}})}.
	\]
	Let $C \subseteq \Sigma^+$ be the set of shortest non-empty
	words $u \in \Sigma^+$ such that $\mu(p_0, u) = p_0$.
	Then $C = \emptyset$ (if $p_0$ has no outgoing transitions)
	or $C$ is a prefix code.
	We then have $L(\mathcal B_{p_0,\{p_0\}}) = C^*$.	
	
	For a direct proof, first note that we have $L(\mathcal B_{p_0,\{p_0\}}) 
	\subseteq \pref{L(\mathcal B)}$
	since for $u \in \Sigma^*$ with $\mu(p_0, u) = p_0$
	we have $\mu(p_0, uv) \in F$
	for some $v \in \Sigma^*$.
	We also have $L(\mathcal B) \subseteq \pref{L(\mathcal B_{p_0,\{p_0\}})}$
	since for $u \in L(\mathcal B)$ we have $\mu(p_0, uv) = p_0$
	for some $v \in \Sigma^*$.
	This gives $\pref{L(\mathcal B)} = \pref{L(\mathcal B_{p_0,\{p_0\}})}$.
	Then proceed as above.
\end{proof}

Note that $C^*$ is sparse if and only if $|C| \le 1$ for a prefix code $C \subseteq \Sigma^+$ by \cref{constr_sync:lem:two-words-code,constr_sync:cor:code_growth_of_star}.

\begin{corollary}
\label{constrsync:cor:SCC}
Let $\automat{B}{P}{\Sigma}{\mu}{p_0}{F}$ be a strongly connected PDFA.
Then the problem $\constrsync{L(\mathcal B)}$ is in \NL,
and we have $L(\mathcal B) = \{\varepsilon\}$ (and so $\constrsync{L(\mathcal B)}$ is solvable in constant time)
or $L(\mathcal B)$ is infinite
and thus $\constrsync{L(\mathcal B)}$ is \LOGSPACE-complete if $L(\mathcal B)$ is also sparse.
\end{corollary}
\begin{proof}
	By \cref{constrsync:thm:SCC-prefix-code}
	we have $\constrsync{L(\mathcal B)}
	\eqACZEROm \constrsync{C^*}$
	for $C = \emptyset$ or some prefix code $C \subseteq \Sigma^+$.
	We have $C = \emptyset$
	precisely if $L(\mathcal B) = \{\varepsilon\}$.
	So, suppose $C \ne \emptyset$ in the following.
	By \cref{thm:reg_monoid} (or \cref{constr_syn:cor:start_equals_final_state}
	since $C^*$ is acceptable by a PDFA with a single final state equal to its start state),
	we find that $\constrsync{L(\mathcal B)}$ is in $\NL$.
	For $C \ne \emptyset$,
	from the method of proof of \cref{constrsync:thm:SCC-prefix-code} and construction of~$C$,
	we find that $\constrsync{L(\mathcal B)}$
	is infinite (we also find that it is sparse if $C = \{u\}$
	for some $u \in \Sigma^*$).	
\end{proof}

Showing \NL-hardness for general regular monoid languages of exponential growth is an \emph{open problem},
even for those of the form~$C^*$ for a prefix code $C$ that is regular and contains at least two words.

\begin{corollary}
	\label{constr_sync:cor:one-state}
	Let $\automat{B}{\{p_0\}}{\Sigma}{\mu}{p_0}{\{p_0\}}$ 
	be a one-state PDFA and $\Gamma = \{ a \in \Sigma \mid \applytrans{\mu}{p_0}{a} = p_0 \}$.
	Then $L(\mathcal B) = \Gamma^*$ and $\constrsync{L(\mathcal B)}$
	has the following complexities:
	\[
	\begin{cases}
		\text{constant time} & \text{if } \Gamma = \emptyset, \\
		\text{\LOGSPACE-complete} & \text{if } |\Gamma| = 1, \\
		\text{\NL-complete} 	& \text{if } |\Gamma| \ge 2.
	\end{cases}
	\]	
\end{corollary}
\begin{proof}
	By \cref{constrsync:rem:alphabet}, 
	we can assume $\Sigma = \Gamma$.
	By \cref{constrsync:cor:SCC}
	the claim follows for $|\Gamma| \le 1$.
	For $|\Gamma| \ge 2$, the claim follows from \crefitem{constr_sync:thm:introductory_examples}{constrsync:ex:item:ab}.		
\end{proof}